\begin{abstract}
    Este informe presenta el diseño, implementación y evaluación de un sistema de visión artificial basado en Deep Learning para clasificar texturas de neumáticos en dos categorías: sano y dañado. Con el fin de mitigar el desbalance de clases severo inherente al conjunto de datos original, se implementaron técnicas complementarias a nivel de datos (mediante \texttt{WeightedRandomSampler}) y a nivel de función de costo (mediante \texttt{Focal Loss}). Se llevó a cabo un estudio de ablación comparando una red neuronal convolucional personalizada diseñada desde cero (\texttt{CustomCNN}) y un modelo preentrenado de gran escala (\texttt{ResNet-50}) mediante Transfer Learning. Los resultados experimentales confirman la superioridad indiscutible de \texttt{ResNet-50}, la cual alcanzó una exactitud de hasta 79.69\% en su configuración de ablación con entropía cruzada y un área bajo la curva ROC (AUC-ROC) de 0.8928 utilizando Focal Loss, superando notablemente al modelo CustomCNN baseline (exactitud de 66.77\% y AUC-ROC de 0.7713). Adicionalmente, se integró el método de interpretabilidad visual Grad-CAM para validar las zonas de activación de la red convolucional y un análisis cualitativo de los principales modos de fallo del sistema para proponer futuras líneas de desarrollo.
\end{abstract}